\documentclass[../main.tex]{subfiles}

\graphicspath{{\subfix{../images/}}}

\begin{document}


\clearpage
\thispagestyle{empty}
\vspace*{7cm}
\section{Introduzione - Parte II}
\vspace*{1.5cm}






\clearpage
\thispagestyle{empty}
\vspace*{7cm}
\section{Introduzione ai sistemi di controllo}
\vspace*{1.5cm}




\subsection{Definizioni utili}

Passiamo ora velocemente a dare una sfilza di definizioni nonché abbreviazioni che saranno utilizzate all'interno di questa sezione.
Notiamo come alcune di esse saranno utilizzate molto di frequente mentre altre verranno solo richiamate. Riteniamo comunque utile riportarle di seguito.


\begin{enumerate}
	\item[\textbf{-}] \textbf{DdB:} questa abbreviazione viene utilizzata per andare ad indicare il \textbf{Diagramma di Bode} di un determinato fattore. In alcuni casi, seppur sbagliando, é riportata tale dicitura seguita da delle parentesi tonde contenenti un carattere o simbolo referente ad un determinato fattore.
	
	\item[\textbf{-}] \textbf{FdT:} questa abbreviazione viene utilizzata per andare ad indicare il \textbf{Funzione di Trasferimento}.
	
	\item[\textbf{-}] \textbf{DdN:} questa abbreviazione viene utilizzata per andare ad indicare il \textbf{Diagramma di Nyquits}.
	
	\item[\textbf{-}] \textbf{BF:}
	
	\item[\textbf{-}] \textbf{AF:}
	
	
\end{enumerate}




\end{document}